\chapter*{Abstract}

Evaluating emotion recognition in text commonly assumes that each text has a single correct
emotion label. Many texts, however, admit several plausible emotional readings, and the one a
given reader settles on depends not only on the words but on who they suppose wrote them. Pride, embarrassment, or regret can be conveyed by the
sentence \emph{``I can't believe I actually did that''} depending on who is taken to be the
speaker. The thesis asks whether perceived author identity measurably affects the interpretation
of emotion, both for human readers and for large language models, and whether the responses of
the two are similar.

Three hundred brief self-reports of emotional experience were sampled from the crowd-enVent
corpus and stratified into three categories of one hundred texts each: \emph{unambiguous control
texts}, \emph{author-independent ambiguous texts}, and \emph{author-relevant ambiguous texts}.
Each text was neutrally framed and also presented under two opposite author personas, with the
text itself left unchanged. Human judgment was gathered across all three framings. Seven large
language models then categorised every text under four framings and six prompt
configurations. The outcome is 7{,}200 model-ensemble distributions and 900 human distributions
(one per text and framing condition), constructed using as many as 50{,}400 model answers and
2{,}899 human ratings.

The hypotheses are tested in two circumstances: one where a single majority or top emotion is
chosen for each condition, and another using the complete \textbf{probability distribution over
eleven emotions}. When the two are compared to each other, one can see what the single-label
method is not able to capture. Among the 300 sampled texts, 19.7\% have no emotion commanding
a strict majority of readings. Of those, 6.3\% still have a unique most-frequent (plurality) label, so a
pipeline that reports only the top label can retain it while covering up how divided readers
really were; the remaining 13.3\% have no single best choice whatsoever, and some arbitrary
tie-break must be made. A distribution preserves that disagreement rather than eliminating it. It
also allows one to ask not only \emph{whether} an answer changed but \emph{how far} and
\emph{in which direction} it changed.

Four hypotheses were tested. As predicted, human judgments shift significantly under
author framing in author-relevant ambiguous texts, while unambiguous texts show no detectable
shift. There is only one persona contrast on which the author-independent category is significant,
and there is no evidence that the two persona effects are different. The model ensemble, however,
shifts significantly under both personas. A previously unanalysed \texttt{generic\_human} control
separates
out two comparisons with matched models: an unspecified author line alone produces about an
8\% move in classifications on both sides, and going further to a named persona is associated
with a separately measured 16--19\% move (the two distances are not additive elements of a single
effect). The predicted category ordering holds in raw shift magnitude. A sensitivity analysis
referring each text to its own conditional sampling null, however, shows that at this sample size
the ordering cannot be told apart from the dispersion in the baseline votes. Lastly, models move
in the same \emph{direction} as human shifts better than chance, though only weakly, and human
shifts are substantially larger in \emph{magnitude} than model shifts. This difference in
magnitude does not vanish when both sides are standardised to exactly the same number of votes.

Two outcomes are reported as negative or near-negative rather than smoothed out. After testing
both human and model shifts, virtually no single per-emotion shift passes multiple-comparison
correction (7 of 66, forming a small but directionally coherent pattern on the model side, rather
than scattered noise). And the exploratory analysis did not detect a prompt-design effect on the size of the
model framing effect. Another finding is methodological in nature, and it reveals most sharply
why the representation matters. Run the same statistical test on the chosen label instead of the
entire distribution, and it rejects the null in \textbf{1} of 6 tests establishing H1 under a
reverse-alphabetical tie-break rule, and \textbf{0} under the alphabetical rule used throughout
this thesis (only these two otherwise-arbitrary rules were tested); the distributional test
rejects it in 3 of them regardless. Using a chosen label, this study would have detected a
statistically significant human framing effect in at best one of six H1 tests, or none at all
under its own convention. More broadly, across all comparisons, two fifths of both human and model
comparisons shifted in ways a selected-label analysis would have recorded as no change whatsoever.
This value is also convention-dependent: when the label method is credited with tracking the full
set of tied top emotions rather than only the alphabetically chosen one, the human figure drops to
about a quarter (25.8\%), while the model figure stays substantially unchanged (38.7\%).

Collectively, these findings support three conclusions: author framing is an actual and
measurable factor affecting the interpretation of emotion, language models are sensitive to it,
and the resemblance between human and model behaviour is partial and directional rather than
close.
